\chapter*{Введение}
\addcontentsline{toc}{chapter}{Введение}

\textbf{Цель работы:} освоить на практике принципы построения, обучения и комплексной оценки сверточных нейронных сетей (Convolutional Neural Networks, CNN) для решения базовых прикладных задач компьютерного зрения, исследовать влияние слоев пакетной нормализации (BatchNormalization) и параметров стохастического градиентного спуска (SGD) на сходимость оптимизационного процесса и обобщающую способность модели.

В соответствии с методическими указаниями к лабораторной работе №~5 кафедры Прикладной математики и кибернетики СибГУТИ, распределение индивидуального задания определяется номером студента в журнале группы: $N_{\text{сп}} = 9$.

\textbf{Параметры индивидуального варианта:}
\begin{enumerate}
    \item \textit{Выбор набора данных:}
    \begin{equation*}
        N_{\text{датасет}} = (N_{\text{сп}} \bmod 6) + 1 = (9 \bmod 6) + 1 = 4 \implies \textbf{Oxford Flowers 102}.
    \end{equation*}
    Коллекция Oxford Flowers 102 содержит 8189 изображений цветковых растений, сгруппированных по 102 классам (от 40 до 258 изображений на категорию). Выборка канонически разделена на обучающую (1020 кадров, строго по 10 примеров на класс), валидационную (1020 кадров) и генеральную тестовую (6149 кадров, более 75\,\% всего объема).
    \item \textit{Выбор архитектуры нейронной сети и оптимизатора:}
    \begin{equation*}
        N_{\text{архитектура}} = (N_{\text{сп}} \bmod 8) + 1 = (9 \bmod 8) + 1 = 2 \implies \textbf{Вариант 2}.
    \end{equation*}
    \begin{itemize}
        \item \textbf{Архитектура CNN:} 3 сверточных слоя (Conv2D) + 2 полносвязных слоя (Dense);
        \item \textbf{Регуляризация:} пакетная нормализация (BatchNormalization);
        \item \textbf{Оптимизатор:} стохастический градиентный спуск (SGD) с коэффициентом момента.
    \end{itemize}
\end{enumerate}

\textbf{Основные шаги выполнения работы:}
\begin{enumerate}
    \item \textbf{Шаг 1. Загрузка и подготовка данных:} исследование структуры выборки Oxford Flowers 102, визуализация репрезентативных категорий, масштабирование изображений к размеру $96 \times 96 \times 3$, нормализация интенсивностей каналов и построение конвейера данных \texttt{tf.data} с потоковой аугментацией.
    \item \textbf{Шаг 2. Создание архитектуры нейронной сети:} проектирование сверточной сети строго по Варианту 2 (3 сверточных блока с фильтрами 32, 64, 128, пакетной нормализацией BatchNormalization, нелинейностями ReLU и пулингом MaxPooling2D, объединенных с 2 полносвязными слоями на 256 и 102 нейрона) и компиляция с оптимизатором SGD.
    \item \textbf{Шаг 3. Обучение, анализ графиков и оценка точности:} обучение модели на 12 эпохах, анализ кривых сходимости функций потерь и метрик точности (Top-1 и Top-5), финальная оценка качества на генеральной тестовой выборке из 6149 изображений, сравнение с уровнем случайного угадывания и качественный анализ предсказаний.
\end{enumerate}

\vspace{12pt}
\noindent \textbf{Ссылка на интерактивный блокнот в Google Colab:} \\
{\footnotesize\url{https://colab.research.google.com/github/MichaelRayven/Machine-Learning-Algorithms/blob/feat/lab5/src/lab5/lab5.ipynb}}

\vspace{6pt}
\noindent \textbf{Репозиторий проекта на GitHub:} \\
{\footnotesize\url{https://github.com/MichaelRayven/Machine-Learning-Algorithms/tree/feat/lab5}}
